% =====================================================================
%  courses/aritmetica/semana_01_01_leyes_logica_propocicional/ejercicios.tex — ejercicios propuestos: leyes de la lógica proposicional
% ---------------------------------------------------------------------
%  Curso: Aritmética | Semana: 01
% =====================================================================

\newpage
\section{Ejercicios propuestos}

\begin{tcolorbox}[
  enhanced, colback=GrisClaro, colframe=AzulPizarra,
  arc=2pt, boxrule=0.7pt, left=8pt, right=8pt, top=4pt, bottom=4pt
]
\small\sffamily
\textbf{Instrucciones:} Marca la alternativa correcta. Cada ejercicio
tiene una única respuesta válida. Aplica las leyes del álgebra de
proposiciones para simplificar antes de evaluar.
\hfill \textbf{Semana 01 --- L\'{o}gica Proposicional}
\end{tcolorbox}
\vspace{0.4cm}

% --- NIVEL BÁSICO ----------------------------------------------------------
\begin{tcolorbox}[
  enhanced, colback=VerdeSuave!40, colframe=VerdeOliva,
  arc=2pt, boxrule=0.5pt, fonttitle=\sffamily\bfseries\small,
  title={Nivel B\'{a}sico}, left=4pt, right=4pt, top=3pt, bottom=3pt
]\end{tcolorbox}

\begin{multicols}{2}
\setcounter{numejercicio}{0}

\ejercicio ¿Cuál es el equivalente más simple de $\sim(\sim p \vee \sim q)$?
\begin{alternativas}
  \item $p \vee q$
  \item $\sim p \wedge \sim q$
  \item $p \wedge q$
  \item $\sim p \vee q$
  \item $p \rightarrow q$
\end{alternativas}
\vspace{0.3cm}

\ejercicio La expresión $p \wedge (p \vee q)$ es equivalente a:
\begin{alternativas}
  \item $p \vee q$
  \item $p \wedge q$
  \item $q$
  \item $p$
  \item $\mathbf{V}$
\end{alternativas}
\vspace{0.3cm}

\ejercicio Según la ley de la condicional, $p \rightarrow q$ es equivalente a:
\begin{alternativas}
  \item $p \wedge \sim q$
  \item $\sim p \wedge q$
  \item $\sim p \vee q$
  \item $p \vee \sim q$
  \item $\sim q \vee p$
\end{alternativas}
\vspace{0.3cm}

\ejercicio Simplificar: $p \vee (p \wedge q)$
\begin{alternativas}
  \item $q$
  \item $p \wedge q$
  \item $\mathbf{V}$
  \item $p \vee q$
  \item $p$
\end{alternativas}
\vspace{0.3cm}

\ejercicio La expresión $p \vee \sim p$ es equivalente a:
\begin{alternativas}
  \item $\mathbf{F}$
  \item $p$
  \item $\mathbf{V}$
  \item $\sim p$
  \item $p \wedge \sim p$
\end{alternativas}
\vspace{0.3cm}

\end{multicols}

% --- NIVEL INTERMEDIO ------------------------------------------------------
\begin{tcolorbox}[
  enhanced, colback=AzulClaro!40, colframe=AzulMedio,
  arc=2pt, boxrule=0.5pt, fonttitle=\sffamily\bfseries\small,
  title={Nivel Intermedio}, left=4pt, right=4pt, top=3pt, bottom=3pt
]\end{tcolorbox}

\begin{multicols}{2}

\ejercicio Simplificar la expresión:
$p \wedge (\sim p \vee q)$
\begin{alternativas}
  \item $p \vee q$
  \item $\sim p \wedge q$
  \item $p \wedge q$
  \item $q$
  \item $p$
\end{alternativas}
\vspace{0.3cm}

\ejercicio La expresión $\sim(p \vee q)$ es equivalente a:
\begin{alternativas}
  \item $\sim p \vee \sim q$
  \item $p \wedge q$
  \item $\sim p \wedge \sim q$
  \item $p \vee \sim q$
  \item $\sim p \rightarrow q$
\end{alternativas}
\vspace{0.3cm}

\ejercicio Reducir: $[(p \rightarrow q) \wedge q] \wedge [(q \rightarrow p) \wedge p]$
\begin{alternativas}
  \item $p \vee q$
  \item $p \wedge q$
  \item $p$
  \item $q$
  \item $\mathbf{V}$
\end{alternativas}
\vspace{0.3cm}

\ejercicio Simplificar la expresión:
$(p \rightarrow q) \wedge (q \rightarrow p)$
\begin{alternativas}
  \item $p \rightarrow q$
  \item $p \vee q$
  \item $p \wedge q$
  \item $p \leftrightarrow q$
  \item $\sim p \vee \sim q$
\end{alternativas}
\vspace{0.3cm}

\ejercicio Reducir: $p \vee (\sim p \wedge q)$
\begin{alternativas}
  \item $p \wedge q$
  \item $\sim p \vee q$
  \item $p \vee q$
  \item $q$
  \item $p$
\end{alternativas}
\vspace{0.3cm}

\end{multicols}

% --- NIVEL AVANZADO --------------------------------------------------------
\begin{tcolorbox}[
  enhanced, colback=NaranjaSuave!60, colframe=NaranjaTierra,
  arc=2pt, boxrule=0.5pt, fonttitle=\sffamily\bfseries\small,
  title={Nivel Avanzado --- Tipo Admisi\'{o}n},
  left=4pt, right=4pt, top=3pt, bottom=3pt
]\end{tcolorbox}

\begin{multicols}{2}

\ejercicio \textbf{(PUCP)} Determinar el equivalente más
simple de:
\[
  (p \vee q) \rightarrow (\sim p \wedge q)
\]
\begin{alternativas}
  \item $q$
  \item $\sim q$
  \item $\sim p$
  \item $p \wedge q$
  \item $\mathbf{V}$
\end{alternativas}
\vspace{0.3cm}

\ejercicio \textbf{(UNI)} Simplificar:
\[
  (p \rightarrow \sim q) \wedge (\sim q \vee \sim p)
\]
\begin{alternativas}
  \item $\sim p \vee \sim q$
  \item $p \wedge q$
  \item $\sim p \wedge q$
  \item $p \rightarrow \sim q$
  \item $\sim p \rightarrow q$
\end{alternativas}
\vspace{0.3cm}

\ejercicio \textbf{(UNMSM)} Reducir la proposición:
``Si el triángulo tiene dos lados iguales, entonces
se llama isósceles; y el triángulo no se llama
isósceles; luego el triángulo no tiene dos lados
iguales.''
\begin{alternativas}
  \item $p \wedge q$
  \item $\mathbf{F}$
  \item $q$
  \item $\mathbf{V}$
  \item $\sim p \vee \sim q$
\end{alternativas}
\vspace{0.3cm}

\ejercicio \textbf{(UNI 2012-I)} Simplificar:
\[
  [(p \rightarrow q) \rightarrow p] \wedge
  [\sim p \rightarrow (\sim p \rightarrow q)]
\]
\begin{alternativas}
  \item $p \vee q$
  \item $p$
  \item $q$
  \item $p \wedge q$
  \item $\mathbf{V}$
\end{alternativas}
\vspace{0.3cm}

\ejercicio Se define la operación
$p \ast q \equiv (p \wedge \sim q) \vee (p \wedge q)$.
Simplificar: $\sim[(p \ast \sim q) \rightarrow (\sim p \ast q)]$
\begin{alternativas}
  \item $\sim p \vee q$
  \item $p \wedge \sim q$
  \item $p \wedge q$
  \item $\sim p \wedge \sim q$
  \item $p$
\end{alternativas}
\vspace{0.3cm}

\end{multicols}

% --- CLAVE DE RESPUESTAS ---------------------------------------------------
\vspace{0.5cm}
\begin{tcolorbox}[
  enhanced, colback=GrisClaro, colframe=GrisMedio,
  arc=2pt, boxrule=0.5pt,
  fonttitle=\sffamily\bfseries\small\color{GrisCarbón},
  title={Clave de Respuestas}
]
\small\sffamily
\begin{multicols}{5}
\noindent
1.~C \quad 2.~D \quad 3.~C \quad 4.~E \quad 5.~C \\
6.~C \quad 7.~C \quad 8.~B \quad 9.~D \quad 10.~C \\
11.~C \quad 12.~A \quad 13.~D \quad 14.~B \quad 15.~B
\end{multicols}
\end{tcolorbox}

% FIN — ejercicios.tex
